% Part: proof-theory
% Chapter: sequent-calculus
% Section: proof-examples

\documentclass[../../../include/open-logic-section]{subfiles}

\begin{document}

\olfileid{pt}{seq}{ex}

\olsection{أمثلة في إنشاء البراهين}

إذا أردنا إنشاء~!!{proof}s لمتتاليات، فأرفق الطرق أن نبدأ بالمتتالية
النهائية ثم نصعد منها: نعين فيها~!!a{formula} نجعلها~!!{formula} الفاعلة،
ونثبت فوقها مقدمات القاعدة التي تلائمها، ثم نعاود هذا الصنيع حتى ننتهي
إلى بديهيات. ولنجعل مطلوبنا، على سبيل المثال، برهنة المتتالية
\[(!C \land !D) \lif !E \Sequent !C \lif (!D \lif !E)\]
نبدأ بـ$!C\lif(!D\lif !E)$؛ فهي من الصورة~$!A\lif !B$ وتقع في يمين
المتتالية. فإذا قابلنا المتتالية كلها بقاعدة~$\RightR{\lif}$ في
\Log{G1c}،
\[\Axiom$ !A, \OLId{greek-variable}{Gamma} \fCenter \OLId{greek-variable}{Delta}, !B$
\RightLabel{\RightR{\lif}}
\UnaryInf$ \OLId{greek-variable}{Gamma} \fCenter \OLId{greek-variable}{Delta}, !A \lif !B $
\DisplayProof\]
بان لنا أن نضع~$\OLId{greek-variable}{Gamma}=\{(!C\land !D)\lif !E\}$، و$\OLId{greek-variable}{Delta}=\emptyset$،
و$!A\ident !C$، و$!B\ident !D\lif !E$. وعلى هذا يجوز أن يكون آخر
البرهان الاستدلال الآتي:
\[\Axiom$ !C, (!C \land !D) \lif !E \fCenter !D \lif !E$
\RightLabel{\RightR{\lif}}
\UnaryInf$(!C \land !D) \lif !E \fCenter !C \lif (!D \lif !E)$
\DisplayProof\]
ثم نعيد العمل نفسه ونجعل~$!D\lif !E$ صيغةً فاعلةً، فنحصل على
\[\Axiom$ !D, !C, (!C \land !D) \lif !E \fCenter !E$
\RightLabel{\RightR{\lif}}
\UnaryInf$!C, (!C \land !D) \lif !E \fCenter !D \lif !E$
\RightLabel{\RightR{\lif}}
\UnaryInf$(!C \land !D) \lif !E \fCenter !C \lif (!D \lif !E)$
\DisplayProof\]
ثم ننظر في~$(!C\land !D)\lif !E$ ونستعمل قاعدة
$\LeftR{\lif}$ في~\Log{G1c}،
\[\Axiom$ \OLId{greek-variable}{Gamma} \fCenter \OLId{greek-variable}{Delta}, !A$
\Axiom$ !B, \OLId{greek-variable}{Gamma} \fCenter \OLId{greek-variable}{Delta}$
\RightLabel{\LeftR{\lif}}
\BinaryInf$ !A \lif !B, \OLId{greek-variable}{Gamma} \fCenter \OLId{greek-variable}{Delta}$
\DisplayProof\]
فينتج لنا
\[\Axiom$!D, !C \fCenter !E, !C \land !D$
\Axiom$!D, !C, !E \fCenter !E$
\RightLabel{\LeftR{\lif}}
\BinaryInf$ !D, !C, (!C \land !D) \lif !E \fCenter !E$
\RightLabel{\RightR{\lif}}
\UnaryInf$!C, (!C \land !D) \lif !E \fCenter !D \lif !E$
\RightLabel{\RightR{\lif}}
\UnaryInf$(!C \land !D) \lif !E \fCenter !C \lif (!D \lif !E)$
\DisplayProof\]
فإن جعلنا~!!{formula}~$!C\land !D$ في المتتالية العليا اليسرى رئيسة،
ثم أعملنا قاعدة~$\RightR{\land}$، تم لنا
\[
\Axiom$!D, !C \fCenter !E, !C$
\Axiom$!D, !C \fCenter !E, !D$
\RightLabel{\RightR{\land}}
\BinaryInf$!D, !C \fCenter !E, !C \land !D$
\Axiom$!D, !C, !E \fCenter !E$
\RightLabel{\LeftR{\lif}}
\BinaryInf$ !D, !C, (!C \land !D) \lif !E \fCenter !E$
\RightLabel{\RightR{\lif}}
\UnaryInf$!C, (!C \land !D) \lif !E \fCenter !D \lif !E$
\RightLabel{\RightR{\lif}}
\UnaryInf$(!C \land !D) \lif !E \fCenter !C \lif (!D \lif !E)$
\DisplayProof
\]
وإلى هنا يستقيم البناء في~\Log{G3c} أيضًا؛ وذلك أن القواعد المستعملة
إلى الآن مشتركة بين~\Log{G1c} و\Log{G3c}. وقد بلغنا جزءًا من~!!a{proof}
في كل متتالية عليا منه~!!{formula} واحدة في اليسار واليمين، غير أن تلك
المتتاليات لم تصر بديهيات بعد.

وأما في~\Log{G1c}، فينبغي أن نبني~!!{proof}s للمتتاليات العليا ابتداءً
من بديهيات على الصورة~$!A\Sequent !A$. وقواعد الإضعاف تكفينا لذلك؛
فالمتتالية العليا في أقصى اليسار يمكن~!!{prove}ها، مثلًا،
$!D,!C\Sequent !E,!C$ كما يأتي:
\[
\Axiom$!C \fCenter !C$
\RightLabel{\LeftR{\Weakening}}
\UnaryInf$!D, !C \fCenter !C$
\RightLabel{\RightR{\Weakening}}
\UnaryInf$!D, !C \fCenter !E, !C$
\]
وبضم هذه الخطوات يكون لنا في~\Log{G1c}~!!a{proof} كامل على الصورة الآتية:
\[
\Axiom$!C \fCenter !C$
\RightLabel{\LeftR{\Weakening}}
\UnaryInf$!D, !C \fCenter !C$
\RightLabel{\RightR{\Weakening}}
\UnaryInf$!D, !C \fCenter !E, !C$
\Axiom$!D \fCenter !D$
\RightLabel{\LeftR{\Weakening}}
\UnaryInf$!D, !C \fCenter !D$
\RightLabel{\RightR{\Weakening}}
\UnaryInf$!D, !C \fCenter !E, !D$
\RightLabel{\RightR{\land}}
\BinaryInf$!D, !C \fCenter !E, !C \land !D$
\Axiom$!E \fCenter !E$
\RightLabel{\LeftR{\Weakening}}
\UnaryInf$!C, !E \fCenter !E$
\RightLabel{\LeftR{\Weakening}}
\UnaryInf$!D, !C, !E \fCenter !E$
\RightLabel{\LeftR{\lif}}
\BinaryInf$ !D, !C, (!C \land !D) \lif !E \fCenter !E$
\RightLabel{\RightR{\lif}}
\UnaryInf$!C, (!C \land !D) \lif !E \fCenter !D \lif !E$
\RightLabel{\RightR{\lif}}
\UnaryInf$(!C \land !D) \lif !E \fCenter !C \lif (!D \lif !E)$
\DisplayProof
\]
لا تتوقف صورة هذا البرهان، ولا سيما ارتفاعه، على ماهية~!!{formula}s
$!C$ و$!D$ و$!E$. فلو أبدلنا بهذه الرموز في~!!{proof} المتقدم أي
!!{formula}s، لبقي الحاصل~!!{proof} في~\Log{G1c}. ومن ثم كان البرهان في
\Log{G1c}~\emph{تخطيطيًا}.

وأما بديهيات~\Log{G3c} فمتتاليات على الصورة~$!A,\OLId{greek-variable}{Gamma}\Sequent
\OLId{greek-variable}{Delta},!A$، بشرط أن تكون~$!A$ ذرية. لذلك قد تبدو
$!D,!C\Sequent !E,!C$ بديهية في~\Log{G3c}، لأن
$\OLId{greek-variable}{Gamma}=\{!D\}$ و$\OLId{greek-variable}{Delta}=\{!E\}$؛ لكنها لا تكون بديهية \emph{على التحقيق}
إلا إن كانت~$!C$ ذرية. فالجزء الذي بنيناه لا يتم برهانًا في~\Log{G3c}
إلا إذا كانت~$!C$ و$!D$ و$!E$~!!{formula}s ذرية.

وعلى هذا لم نثبت، إذا دققنا، أن
\[\Log{G3c} \Proves (!C \land !D) \lif !E \Sequent !C \lif (!D \lif
!E),\]
غير أن هذا القدر هو المطلوب عمليًا، وهو غاية ما يتيحه الرسم التخطيطي؛
وذلك أن كل متتالية على الصورة~$!A,\OLId{greek-variable}{Gamma}\Sequent\OLId{greek-variable}{Delta},!A$ قابلة
للبرهنة في~\Log{G3c}، وإن لم تكن~$!A$ ذرية. وبرهان ذلك بالاستقراء على
عمق~$\depth{!A}$ للصيغة~$!A$.

\begin{defn}\ollabel{defn:depth}
نجعل \emph{عمق}~$\depth{!A}$ لـ~!!a{formula} معرفًا على الوجه الآتي:
\begin{enumerate}
  \item إذا كانت~$!A$ ذرية، فإن~$\depth{!A}=\OLMathNumeral{0}$.
  \item إذا كانت~$!A\ident\lnot !B$، أو
    $!A\ident\lforall[\OLId{latin-ordinary}{x}][!B(\OLId{latin-ordinary}{x})]$، أو
    $!A\ident\lexists[\OLId{latin-ordinary}{x}][!B(\OLId{latin-ordinary}{x})]$، فإن
    $\depth{!A}=\depth{!B}+\OLMathNumeral{1}$.
  \item إذا كانت~$!A\ident(!B\land !C)$، أو
    $!A\ident(!B\lor !C)$، أو $!A\ident(!B\lif !C)$، فإن
    $\depth{!A}=\max(\depth{!B},\depth{!C})+\OLMathNumeral{1}$.
\end{enumerate}
\end{defn}

ومن اليسير أن نثبت~$\depth{!A(\OLId{latin-ordinary}{x})}=\depth{!A(\OLId{latin-ordinary}{t})}$ لكل حد~$\OLId{latin-ordinary}{t}$. والذي يعنينا
أن عمق~!!{formula} الرئيسة في القاعدة يفوق أبدًا عمق
!!{formula}s الفاعلة؛ ومن أمثلة ذلك:
\begin{align*}
\depth{\OLId{latin-looped}{P}(\OLId{latin-ordinary}{x})} = \depth{\OLId{latin-looped}{P}(\OLId{latin-ordinary}{a})} & = \OLMathNumeral{0},\\
\depth{\lforall[\OLId{latin-ordinary}{x}][\OLId{latin-looped}{P}(\OLId{latin-ordinary}{x})]} & = \OLMathNumeral{1},\\
\depth{(\lforall[\OLId{latin-ordinary}{x}][\OLId{latin-looped}{P}(\OLId{latin-ordinary}{x})] \lif \OLId{latin-looped}{P}(\OLId{latin-ordinary}{a}))} & = \max(\OLMathNumeral{1},\OLMathNumeral{0}) + \OLMathNumeral{1} = \OLMathNumeral{2}.
\end{align*}
وبهذا المقياس نستطيع إثبات نتائج بالاستقراء على~$\depth{!A}$، ومنها ما يأتي.

\begin{prop}\ollabel{prop:G1c-axioms}
لكل~$!A$ يوجد في~$\Log{G1c}$~!!a{proof} لـ~$!A\Sequent !A$، وجميع
بديهيات ذلك البرهان ذرية.
\end{prop}

\begin{proof}
  نستقرئ على~$\depth{!A}$. فإن كان~$\depth{!A}=\OLMathNumeral{0}$، كانت~$!A$ ذرية،
  وكانت~$!A\Sequent !A$ بذاتها~!!{proof} في~\Log{G1c} على الوجه
  المطلوب. وإن كان~$\depth{!A}>\OLMathNumeral{0}$، كانت~$!A$ مؤلفة من~!!{formula}s
  دونها عمقًا؛ كأن تكون~$!A\ident\lnot !B$، أو~$!A\ident(!B\land !C)$،
  أو~$!A\ident\lforall[\OLId{latin-ordinary}{x}][!B(\OLId{latin-ordinary}{x})]$. وفرض الاستقراء أنه لكل~$!D$ يحقق
  $\depth{!D}<\depth{!A}$، لدينا
  $\Log{G1c}\Proves !D\Sequent !D$ انطلاقًا من بديهيات ذرية.

  فأما إن كانت~$!A\ident\lnot !B$، فإن~$\depth{!B}<\depth{!A}$، ومن فرض
  الاستقراء نحصل على~!!a{proof}~$\pi_\OLMathNumeral{1}$ في~\Log{G1c} بديهياته ذرية
  لـ~$!B\Sequent !B$. ثم نمده إلى~!!a{proof} لـ
  $\lnot !B\Sequent\lnot !B$ كما يأتي:
  \[
  \AxiomC{}
  \RightLabel{$\pi_\OLMathNumeral{1}$}
  \Deduce$!B \fCenter !B$
  \RightLabel{\LeftR{\lnot}}
  \UnaryInf$\lnot !B, !B \fCenter $
  \RightLabel{\RightR{\lnot}}
  \UnaryInf$\lnot !B \fCenter \lnot !B$
  \DisplayProof
\]

  وأما إن كانت~$!A\ident(!B\land !C)$، فإن~$\depth{!B}<\depth{!A}$ و
  $\depth{!C}<\depth{!A}$. وبفرض الاستقراء يوجد في
  \Log{G1c}~!!{proof}s~$\pi_\OLMathNumeral{1}$ و$\pi_\OLMathNumeral{2}$ بديهياتها ذرية لـ
  $!B\Sequent !B$ و$!C\Sequent !C$. ومنهما نركب~!!a{proof}
  لـ~$!B\land !C\Sequent !B\land !C$ كما يأتي:
  \[
  \AxiomC{}
  \RightLabel{$\pi_\OLMathNumeral{1}$}
  \Deduce$!B \fCenter !B$
  \RightLabel{\LeftR{\Weakening}}
  \UnaryInf$!B, !C \fCenter !B$
  \RightLabel{\LeftR{\land}}
  \UnaryInf$!B \land !C, !C \fCenter !B$
  \RightLabel{\LeftR{\land}}
  \UnaryInf$!B \land !C, !B \land !C \fCenter !B$
  \AxiomC{}
  \RightLabel{$\pi_\OLMathNumeral{2}$}
  \Deduce$!C \fCenter !C$
  \RightLabel{\LeftR{\Weakening}}
  \UnaryInf$!B, !C \fCenter !C$
  \RightLabel{\LeftR{\land}}
  \UnaryInf$!B \land !C, !C \fCenter !C$
  \RightLabel{\LeftR{\land}}
  \UnaryInf$!B \land !C, !B \land !C \fCenter !C$
  \RightLabel{\RightR{\land}}
  \BinaryInf$!B \land !C, !B \land !C \fCenter !B \land !C$
  \RightLabel{\LeftR{\Contraction}}
  \UnaryInf$!B \land !C \fCenter !B \land !C$
  \DisplayProof
  \]

  وأما إذا كانت~$!A\ident\lforall[\OLId{latin-ordinary}{x}][!B(\OLId{latin-ordinary}{x})]$، فخذ~!!a{constant}~$\OLId{latin-ordinary}{c}$
  لا يظهر في~$\lforall[\OLId{latin-ordinary}{x}][!B(\OLId{latin-ordinary}{x})]$. وحينئذ
  $\depth{!B(\OLId{latin-ordinary}{c})}=\depth{!B(\OLId{latin-ordinary}{x})}<\depth{!A}$، وبفرض الاستقراء يوجد
  !!a{proof}~$\pi_\OLMathNumeral{1}$ في~\Log{G1c} بديهياته ذرية لـ
  $!B(\OLId{latin-ordinary}{c})\Sequent !B(\OLId{latin-ordinary}{c})$. ومنه ننشئ~!!a{proof} لـ
  $\lforall[\OLId{latin-ordinary}{x}][!B(\OLId{latin-ordinary}{x})]\Sequent\lforall[\OLId{latin-ordinary}{x}][!B(\OLId{latin-ordinary}{x})]$ كما يأتي:
  \[
  \AxiomC{}
  \RightLabel{$\pi_\OLMathNumeral{1}$}
  \Deduce$!B(\OLId{latin-ordinary}{c}) \fCenter !B(\OLId{latin-ordinary}{c})$
  \RightLabel{\LeftR{\lforall}}
  \UnaryInf$\lforall[\OLId{latin-ordinary}{x}][!B(\OLId{latin-ordinary}{x})] \fCenter !B(\OLId{latin-ordinary}{c})$
  \RightLabel{\RightR{\lforall}}
  \UnaryInf$\lforall[\OLId{latin-ordinary}{x}][!B(\OLId{latin-ordinary}{x})] \fCenter \lforall[\OLId{latin-ordinary}{x}][!B(\OLId{latin-ordinary}{x})]$
  \DisplayProof
\]

  وأما سائر الحالات فتُستكمل تمرينًا.
\end{proof}

\begin{prob}
  أتمم برهان~\olref{prop:G1c-axioms} بمعالجة الحالات التي يكون فيها
  $!A\ident(!B\lor !C)$، و$!A\ident(!B\lif !C)$، و
  $!A\ident\lexists[\OLId{latin-ordinary}{x}][!B(\OLId{latin-ordinary}{x})]$.
\end{prob}

وكان يجزئنا في هذا الموضع أيضًا استعمال~!!{proof} الآتي
\[
\AxiomC{}
\RightLabel{$\pi_\OLMathNumeral{1}$}
\Deduce$!B \fCenter !B$
\RightLabel{\RightR{\lnot}}
\UnaryInf$\fCenter !B, \lnot !B$
\RightLabel{\LeftR{\lnot}}
\UnaryInf$\lnot !B \fCenter \lnot !B$
\DisplayProof
\]
في حالة~$!A\ident\lnot !B$. إلا أن هذا الوجه لا يصح في بعض حسابات
المتتاليات؛ فإن حساب المتتاليات الحدسي~\Log{G1i} يماثل~\Log{G1c}، غير
أنه لا يجيز في يمين المتتالية أكثر من~!!{formula} واحدة. فإذا كانت
$\OLId{greek-variable}{Delta}=\emptyset$، كان~!!{proof} الأول~!!a{proof} في~\Log{G1i} كذلك؛
وأما الثاني فلا، لاجتماع~$!B$ و$\lnot !B$ في يمين إحدى متتالياته.

ولا يجوز، حتى في~\Log{G1c}، أن نعكس ترتيب القاعدتين في حالة
$!A\ident\lforall[\OLId{latin-ordinary}{x}][!B(\OLId{latin-ordinary}{x})]$؛ فالشجرة الآتية \emph{ليست}~!!a{proof}:
  \[
  \AxiomC{}
  \RightLabel{$\pi_\OLMathNumeral{1}$}
  \Deduce$!B(\OLId{latin-ordinary}{c}) \fCenter !B(\OLId{latin-ordinary}{c})$
  \RightLabel{\RightR{\lforall}}
  \UnaryInf$!B(\OLId{latin-ordinary}{c}) \fCenter \lforall[\OLId{latin-ordinary}{x}][!B(\OLId{latin-ordinary}{x})]$
  \RightLabel{\LeftR{\lforall}}
  \UnaryInf$\lforall[\OLId{latin-ordinary}{x}][!B(\OLId{latin-ordinary}{x})] \fCenter \lforall[\OLId{latin-ordinary}{x}][!B(\OLId{latin-ordinary}{x})]$
  \DisplayProof
\]
إذ تنقض شرط المتغير المميز في~$\RightR{\lforall}$.

\begin{prop}\ollabel{prop:G3c-axioms}
كل متتالية على الصورة~$!A,\OLId{greek-variable}{Gamma}\Sequent\OLId{greek-variable}{Delta},!A$ قابلة للبرهنة
في~\Log{G3c}، وإن كانت~$!A$ \emph{غير مشترط} فيها أن تكون ذرية.
\end{prop}

\begin{proof}
هذه الحجة قريبة من برهان~\olref{prop:G1c-axioms}، غير أن ضبط فرض الاستقراء
هنا لازم. فأما حالة~$\OLId{latin-ordinary}{n}=\depth{!A}=\OLMathNumeral{0}$ فبينة: إذا كان~$\OLId{latin-ordinary}{n}=\OLMathNumeral{0}$، كانت~$!A$
ذرية، ومن ثم كانت~$!A,\OLId{greek-variable}{Gamma}\Sequent\OLId{greek-variable}{Delta},!A$ بديهية في~\Log{G3c}.

ولنفرض الآن~$\OLId{latin-ordinary}{n}>\OLMathNumeral{0}$، ثم نفصل الحالات. وفرض الاستقراء أن لكل~$!D$
يحقق~$\depth{!D}<\OLId{latin-ordinary}{n}$، ولكل~$\Pi$ و$\Lambda$، يثبت
$\Log{G3c}\Proves !D,\Pi\Sequent\Lambda,!D$. فإذا كانت الحالة
$!A\ident(!B\land !C)$، استعملنا الفرض مرتين: مرة مع~$!D=!B$،
و$\Pi=!C,\OLId{greek-variable}{Gamma}$، و$\Lambda=\OLId{greek-variable}{Delta}$، فنحصل على~!!a{proof}~$\pi_\OLMathNumeral{1}$ لـ
$!B,!C,\OLId{greek-variable}{Gamma}\Sequent\OLId{greek-variable}{Delta},!B$؛ ومرة مع~$!D=!C$،
و$\Pi=!B,\OLId{greek-variable}{Gamma}$، و$\Lambda=\OLId{greek-variable}{Delta}$، فنحصل على~!!a{proof}~$\pi_\OLMathNumeral{2}$ لـ
$!B,!C,\OLId{greek-variable}{Gamma}\Sequent\OLId{greek-variable}{Delta},!C$. ومنهما يتألف~!!{proof} لـ
$!B\land !C,\OLId{greek-variable}{Gamma}\Sequent\OLId{greek-variable}{Delta},!B\land !C$:
\[
\AxiomC{}
\RightLabel{$\pi_\OLMathNumeral{1}$}
\Deduce$!B, !C, \OLId{greek-variable}{Gamma} \fCenter \OLId{greek-variable}{Delta}, !B$
\RightLabel{\LeftR{\land}}
\UnaryInf$!B \land !C, \OLId{greek-variable}{Gamma} \fCenter \OLId{greek-variable}{Delta}, !B$
\AxiomC{}
\RightLabel{$\pi_\OLMathNumeral{2}$}
\Deduce$!B, !C, \OLId{greek-variable}{Gamma} \fCenter \OLId{greek-variable}{Delta}, !C$
\RightLabel{\LeftR{\land}}
\UnaryInf$!B \land !C, \OLId{greek-variable}{Gamma} \fCenter \OLId{greek-variable}{Delta}, !C$
\RightLabel{\RightR{\land}}
\BinaryInf$!B \land !C, \OLId{greek-variable}{Gamma} \fCenter \OLId{greek-variable}{Delta}, !B \land !C$
\DisplayProof
\]

وأما حالة~$!A\ident\lforall[\OLId{latin-ordinary}{x}][!B(\OLId{latin-ordinary}{x})]$، فاختر~!!a{constant}~$\OLId{latin-ordinary}{c}$ لا يظهر
في~$\lforall[\OLId{latin-ordinary}{x}][!B(\OLId{latin-ordinary}{x})]$ ولا~$\OLId{greek-variable}{Gamma}$ ولا~$\OLId{greek-variable}{Delta}$. ثم أجر فرض الاستقراء
مع~$!D=!B(\OLId{latin-ordinary}{c})$، و$\Pi=\lforall[\OLId{latin-ordinary}{x}][!B(\OLId{latin-ordinary}{x})],\OLId{greek-variable}{Gamma}$، و$\Lambda=\OLId{greek-variable}{Delta}$؛
فيعطينا~!!a{proof}~$\pi_\OLMathNumeral{1}$ لـ
$!B(\OLId{latin-ordinary}{c}),\lforall[\OLId{latin-ordinary}{x}][!B(\OLId{latin-ordinary}{x})],\OLId{greek-variable}{Gamma}\Sequent\OLId{greek-variable}{Delta},!B(\OLId{latin-ordinary}{c})$، ومنه نبني:
\[
\AxiomC{}
\RightLabel{$\pi_\OLMathNumeral{1}$}
\Deduce$!B(\OLId{latin-ordinary}{c}), \lforall[\OLId{latin-ordinary}{x}][!B(\OLId{latin-ordinary}{x})], \OLId{greek-variable}{Gamma} \fCenter \OLId{greek-variable}{Delta}, !B(\OLId{latin-ordinary}{c})$
\RightLabel{\LeftR{\lforall}}
\UnaryInf$\lforall[\OLId{latin-ordinary}{x}][!B(\OLId{latin-ordinary}{x})], \OLId{greek-variable}{Gamma} \fCenter \OLId{greek-variable}{Delta}, !B(\OLId{latin-ordinary}{c})$
\RightLabel{\RightR{\lforall}}
\UnaryInf$\lforall[\OLId{latin-ordinary}{x}][!B(\OLId{latin-ordinary}{x})], \OLId{greek-variable}{Gamma} \fCenter \OLId{greek-variable}{Delta}, \lforall[\OLId{latin-ordinary}{x}][!B(\OLId{latin-ordinary}{x})]$
\DisplayProof
\]
وشرط المتغير المميز في~$\RightR{\lforall}$ متحقق بما اشترطناه في اختيار
$\OLId{latin-ordinary}{c}$.
\end{proof}

\begin{prob}
  أتمم برهان~\olref{prop:G3c-axioms} بمعالجة الحالتين اللتين يكون فيهما
  $!A\ident(!B\lif !C)$ و$!A\ident\lexists[\OLId{latin-ordinary}{x}][!B(\OLId{latin-ordinary}{x})]$.
\end{prob}

\end{document}
